\chapter{Síntese}

\begin{objetivo}
Responder à pergunta com proposições numeradas, cada uma com a sua base
textual, separando o que é claro do que fica em aberto.
\end{objetivo}

Duas proposições de partida, que só descrevem o que os textos dizem
diretamente. Revise-as e acrescente as suas.

\proposicao{Na Lei, o sangue é associado à vida (\hb{נֶפֶשׁ}) e reservado ao
altar; por isso não deveria ser comido.}{Lv 17:10–14; Dt 12:23–25; cf. Gn 9:4}

\proposicao{O Novo Testamento usa a linguagem do sangue para falar do
significado da morte de Jesus.}{Mt 26:28; Ef 1:7; Hb 9:12–14; 1Pe 1:18–19}

\campovazio{Outras proposições}{uma por linha, com os textos que as sustentam}
\campovazio{Resposta à pergunta}{em um parágrafo}
\campovazio{O que fica em aberto}{}
\campovazio{Comparação com o que eu pensava}{(Etapa 1)}
\campovazio{Aplicação}{}
